\chapter{Modulation and I/Q Signals}

So far every signal we have studied has lived at its ``natural'' place on the
frequency axis. A $50\,$Hz hum sits at $50\,$Hz. A muscle twitch smears its energy
across a few tens of hertz near zero. Speech crowds into the band below a few
kilohertz. The Fourier transform told us \emph{where} a signal lives; filtering let
us keep or discard pieces of that real estate. This chapter is about the verb that
acts on frequency itself: \emph{modulation} -- the deliberate act of \emph{moving} a
signal's spectrum from one place to another, and later moving it back.

That single idea -- ``multiply by a sinusoid and the spectrum slides'' -- is the
beating heart of all radio, of every Wi-Fi packet and every FM station, of radar and
sonar and ultrasound, of the lock-in amplifier on a lab bench, and, as we will see,
of the very accelerometer and gyroscope chips that will sit on our wrist tracker. It
is also the doorway to the single most important data structure in modern signal
processing: the \emph{I/Q} (in-phase / quadrature) pair, a complex number per sample
that carries both the amplitude and the phase of a signal relative to a reference
oscillator. Learn to see I/Q geometrically -- as a little arrow spinning in the
complex plane, slowed down until it crawls -- and an enormous amount of ``hard''
engineering becomes obvious.

We build the chapter from the modulation property of the Fourier transform, watch
spectra shift in pictures, then assemble amplitude modulation (AM) and frequency
modulation (FM) as two ways of writing information onto a carrier. We meet the
analytic signal as a preview (the next chapter, on the Hilbert transform, makes it
rigorous), define the complex baseband and the I/Q representation, explain why I/Q is
literally everywhere in radio and sensing hardware, and finish with the basic recipes
for getting the message back out: coherent and envelope demodulation.

\section{The one idea: multiplying by a sinusoid shifts the spectrum}
\label{ch23:mod-property}

Recall the modulation (frequency-shift) property of the Fourier transform, the dual of
the time-shift property. If $x(t) \overset{\FT}{\longleftrightarrow} X(\jj\omega)$, then
multiplying $x$ by a complex exponential shifts its whole spectrum:
\begin{equation}
\label{ch23:expshift}
e^{\jj\omega_c t}\, x(t)
\;\overset{\FT}{\longleftrightarrow}\;
X\!\big(\jj(\omega-\omega_c)\big).
\end{equation}
The proof is one line -- substitute into the transform integral and absorb the
exponential into the kernel:
\begin{equation}
\int_{-\infty}^{\infty} \big[e^{\jj\omega_c t}x(t)\big]\,e^{-\jj\omega t}\,\dd t
= \int_{-\infty}^{\infty} x(t)\,e^{-\jj(\omega-\omega_c)t}\,\dd t
= X\!\big(\jj(\omega-\omega_c)\big).
\end{equation}
Multiplying in time by $e^{\jj\omega_c t}$ \emph{translates} the spectrum to the right
by $\omega_c$. That is the whole game.

\begin{intuitionbox}
Frequency is rotation rate (Chapter on sinusoids and phasors). The signal $x(t)$ is a
recipe of little spinning phasors at various rates. Multiplying every one of them by
$e^{\jj\omega_c t}$ -- itself a phasor spinning at rate $\omega_c$ -- speeds each one up
by exactly $\omega_c$. A phasor that was spinning at $\omega$ now spins at
$\omega+\omega_c$. So the entire collection of frequencies slides up the axis by
$\omega_c$, rigidly, without changing shape. \emph{Modulation is just adding a constant
to everybody's rotation rate.}
\end{intuitionbox}

Real signals do not come with a single complex exponential; they come with cosines.
A cosine is two phasors, one spinning each way, $\cos(\omega_c t) = \tfrac12 e^{\jj\omega_c
t} + \tfrac12 e^{-\jj\omega_c t}$. So multiplying by a real cosine makes \emph{two}
shifted copies, one up and one down:
\begin{equation}
\label{ch23:cosmod}
\cos(\omega_c t)\,x(t)
\;\overset{\FT}{\longleftrightarrow}\;
\tfrac12 X\!\big(\jj(\omega-\omega_c)\big) + \tfrac12 X\!\big(\jj(\omega+\omega_c)\big).
\end{equation}
This is the most-used equation in all of communications. A baseband message $x(t)$
crowded around $\omega = 0$ becomes, after multiplication by a carrier cosine, two
half-height replicas centered at $\pm\omega_c$. The picture is worth a thousand
derivations.

\begin{figure}[H]
\centering
\begin{tikzpicture}
\begin{axis}[width=13cm,height=4.3cm,
  xlabel={frequency $\omega$}, ylabel={$|\cdot|$},
  xmin=-9,xmax=9,ymin=0,ymax=1.25,
  ytick=\empty, xtick={-6,0,6}, xticklabels={$-\omega_c$,$0$,$+\omega_c$},
  axis lines=left, tick align=outside,
  title={Multiplying by $\cos(\omega_c t)$: the baseband triangle splits into two half-height copies}]
% baseband message centered at 0 (dashed, the "before")
\addplot[sigblue,thick,dashed,domain=-2:2,samples=41] {1-abs(x)/2};
% two shifted half-height copies (the "after")
\addplot[specred,thick,domain=4:8,samples=41] {0.5*(1-abs(x-6)/2)};
\addplot[specred,thick,domain=-8:-4,samples=41] {0.5*(1-abs(x+6)/2)};
\node[sigblue] at (axis cs:0,1.13) {\footnotesize $X(\jj\omega)$};
\node[specred] at (axis cs:6,0.66) {\footnotesize copies};
\end{axis}
\end{tikzpicture}
\caption{The modulation property in pictures. A baseband spectrum $X(\jj\omega)$ (blue,
dashed) multiplied by a carrier cosine produces two half-height, frequency-shifted
replicas at $\pm\omega_c$ (red). Demodulation will be the act of sliding one copy back
to the origin.}
\label{ch23:fig-split}
\end{figure}

\begin{whybox}
Why does multiplication in time shift in frequency, while \emph{convolution} in time
multiplies in frequency? They are the same theorem read two ways. Multiplying $x(t)$ by
$\cos\omega_c t$ in time is \emph{convolving} $X(\jj\omega)$ with the spectrum of the
cosine, which is a pair of impulses at $\pm\omega_c$. Convolving any spectrum with an
impulse at $\omega_c$ simply copies it, shifted, to $\omega_c$. Two impulses, two copies.
Everything in this chapter is the convolution theorem wearing different clothes.
\end{whybox}

\begin{connectionbox}
The frequency-shift property is the Fourier dual of the time-shift property from the
properties chapter. Both come from the same algebraic fact: $e^{\jj a t}$ is an
\emph{eigen-thing} of shifting. In the linear-algebra companion this is the statement
that the Fourier basis diagonalizes the shift operator; here it tells us that a pure
tone is the natural ``carrier'' for moving information around the frequency axis.
\end{connectionbox}

\section{Amplitude modulation: writing on the height of a carrier}
\label{ch23:am}

The oldest way to send a message on a radio wave is to let the message control the
\emph{strength} of a fast oscillation. Pick a carrier frequency $\omega_c$ far above
anything in your message $m(t)$ (think $\omega_c$ at $1\,$MHz, $m$ at audio rates below
$5\,$kHz). \emph{Double-sideband suppressed-carrier} AM is the bare product
\begin{equation}
\label{ch23:dsbsc}
s_{\text{DSB}}(t) = m(t)\,\cos(\omega_c t),
\end{equation}
whose spectrum, by \eqref{ch23:cosmod}, is two copies of $M(\jj\omega)$ at $\pm\omega_c$.
It is spectrally efficient but a receiver cannot find the carrier phase from the signal
alone, because when $m(t)$ goes negative the carrier flips phase by $180^\circ$.

Classic broadcast AM fixes this by adding a constant so the envelope never crosses zero:
\begin{equation}
\label{ch23:am-eq}
s_{\text{AM}}(t) = A\big[1 + \mu\, m(t)\big]\cos(\omega_c t),
\qquad \abs{m(t)}\le 1,\ \ 0<\mu\le 1.
\end{equation}
Here $\mu$ is the \emph{modulation index}. As long as $1+\mu m(t) > 0$, the slow
\emph{envelope} $A[1+\mu m(t)]$ traces the message exactly, riding on the fast carrier.
A receiver can then recover $m$ by the crudest possible circuit -- a diode and a
capacitor that follow the peaks. That is why AM radio was buildable in 1920 with a
crystal and a coil.

\begin{figure}[H]
\centering
\begin{tikzpicture}
\begin{axis}[width=13cm,height=4.6cm,
  xlabel={time $t$}, ylabel={amplitude},
  xmin=0,xmax=6.3,ymin=-2.1,ymax=2.1,
  ytick=\empty, xtick=\empty, axis lines=left,
  title={AM: a fast carrier whose height is shaped by the slow message envelope}]
% carrier modulated by 1 + 0.7 cos(x)
\addplot[sigblue,thick,domain=0:6.3,samples=200]
   {(1+0.7*cos(deg(x)))*cos(deg(13*x))};
% upper envelope
\addplot[specred,thick,domain=0:6.3,samples=120] {1+0.7*cos(deg(x))};
% lower envelope
\addplot[specred,thick,dashed,domain=0:6.3,samples=120] {-(1+0.7*cos(deg(x)))};
\end{axis}
\end{tikzpicture}
\caption{Amplitude modulation with $\mu=0.7$. The blue waveform is the transmitted
$s_{\text{AM}}(t)$; the red curve is the envelope $A[1+\mu m(t)]$ that a diode detector
follows. Because the envelope stays positive, the message is read off directly.}
\label{ch23:fig-am}
\end{figure}

\begin{example}[Sidebands of a single tone]
\label{ch23:ex-tone}
Let the message be a pure tone $m(t)=\cos(\omega_m t)$ with $\mu=1$. Then
\[
s_{\text{AM}}(t) = A\cos(\omega_c t) + \tfrac{A}{2}\cos\!\big((\omega_c+\omega_m)t\big)
 + \tfrac{A}{2}\cos\!\big((\omega_c-\omega_m)t\big).
\]
The spectrum is three pairs of lines: the \emph{carrier} at $\pm\omega_c$ plus an
\emph{upper sideband} at $\pm(\omega_c+\omega_m)$ and a \emph{lower sideband} at
$\pm(\omega_c-\omega_m)$. All the \emph{information} lives in the sidebands; the carrier
line carries none of it (it is a constant) yet consumes most of the power. That waste is
exactly what DSB-SC and single-sideband (SSB) schemes attack by removing the carrier,
and one sideband too.
\end{example}

\begin{trapbox}
\textbf{Overmodulation.} If $\mu>1$, the quantity $1+\mu m(t)$ dips below zero and the
envelope ``folds.'' A simple envelope detector then reproduces $\abs{1+\mu m(t)}$, a
rectified, badly distorted version of the message. Keeping $\mu\le 1$ is not a stylistic
choice; it is the condition under which envelope detection is even valid. The same
trap reappears in sensing: an amplitude-modulated vibration read by a peak-follower must
not flip sign, or the recovered envelope is garbage.
\end{trapbox}

\begin{applicationbox}
\textbf{Where AM-style shifting hides.} Beyond radio, the multiply-by-a-carrier trick is
the basis of the \emph{lock-in amplifier}: to measure a tiny signal buried in noise, you
deliberately modulate it at a known frequency $\omega_c$ (chop a light beam, drive a
bridge), then multiply the noisy measurement by $\cos\omega_c t$ and low-pass filter.
Only energy that was sitting at $\omega_c$ survives the trip back to baseband; broadband
noise is shifted but mostly filtered away. Ultrasound imaging, optical heterodyne
detection, and capacitive touch sensing all run on the same principle.
\end{applicationbox}

\section{Frequency modulation: writing on the rotation rate}
\label{ch23:fm}

AM writes the message onto the carrier's \emph{height}. FM writes it onto the carrier's
\emph{rate of rotation}. This is the natural move once you have internalized that
frequency \emph{is} rotation rate: instead of nudging the length of the spinning arrow,
nudge how fast it spins.

The clean way to say ``how fast it is spinning right now'' is the \emph{instantaneous
frequency}. Write any narrowband signal as $s(t)=A\cos\theta(t)$ with a time-varying
phase $\theta(t)$. The instantaneous angular frequency is the rate the phase advances:
\begin{equation}
\label{ch23:instfreq}
\omega_i(t) \;=\; \frac{\dd \theta}{\dd t}.
\end{equation}
For an ordinary tone $\theta(t)=\omega_c t$ and $\omega_i=\omega_c$, a constant -- the arrow
spins at a steady rate. FM makes the rate wobble in step with the message:
\begin{equation}
\label{ch23:fmdef}
\omega_i(t) = \omega_c + k_f\, m(t)
\quad\Longrightarrow\quad
s_{\text{FM}}(t) = A\cos\!\Big(\omega_c t + k_f\!\int_0^t m(\tau)\,\dd\tau\Big).
\end{equation}
The integral appears because phase is the accumulation of frequency; this is the same
``angle is the running total of angular velocity'' relation from the calculus companion.
The constant $k_f$ sets how many radians-per-second of frequency swing you get per unit
of message.

\begin{figure}[H]
\centering
\begin{tikzpicture}
\begin{axis}[width=13cm,height=4.4cm,
  xlabel={time $t$}, ylabel={amplitude},
  xmin=0,xmax=6.3,ymin=-1.5,ymax=1.5,
  ytick=\empty, xtick=\empty, axis lines=left,
  title={FM: constant height, but the spacing of the wiggles breathes with the message}]
% FM: cos(wc t + beta sin(wm t)), here approximate instantaneous freq variation
\addplot[sigblue,thick,domain=0:6.3,samples=200]
   {cos(deg(9*x + 4*sin(deg(x))))};
\end{axis}
\end{tikzpicture}
\caption{Frequency modulation. The amplitude is rock-steady; the \emph{information} is in
the changing spacing of the zero crossings -- the wiggles bunch up where the
instantaneous frequency is high and spread out where it is low.}
\label{ch23:fig-fm}
\end{figure}

The spectrum of FM is genuinely more complicated than AM. Even a single-tone message
$m(t)=\cos\omega_m t$ produces \emph{infinitely many} sidebands, spaced $\omega_m$ apart,
whose heights are Bessel functions $J_n(\beta)$ of the modulation index
$\beta = k_f A_m/\omega_m$. The useful, memorable summary is \emph{Carson's rule}: the
practical bandwidth of an FM signal is
\begin{equation}
\label{ch23:carson}
B \;\approx\; 2(\Delta f + f_m),
\end{equation}
where $\Delta f$ is the peak frequency deviation and $f_m$ the highest message frequency.
Wider deviation buys more noise immunity at the cost of more bandwidth -- the trade FM
broadcasting makes deliberately, which is why FM sounds cleaner than AM but eats more
spectrum.

\begin{figure}[H]
\centering
\begin{tikzpicture}
\begin{axis}[width=13cm,height=4.4cm,
  xlabel={frequency offset from carrier (units of $\omega_m$)}, ylabel={$|J_n(\beta)|$},
  xmin=-5.4,xmax=5.4,ymin=0,ymax=0.72,
  xtick={-4,-3,-2,-1,0,1,2,3,4}, ytick=\empty,
  axis lines=left, tick align=outside,
  title={FM single-tone spectrum at $\beta=2$: a comb of Bessel-weighted sidebands}]
% heights from J_n(2): J0=0.224, J1=0.577, J2=0.353, J3=0.129, J4=0.034
\addplot+[ycomb,sigblue,very thick,mark=*,mark size=1.4pt,samples at={0}] {0.224};
\addplot+[ycomb,specred,very thick,mark=*,mark size=1.4pt,samples at={-1,1}] {0.577};
\addplot+[ycomb,specred,very thick,mark=*,mark size=1.4pt,samples at={-2,2}] {0.353};
\addplot+[ycomb,specred,very thick,mark=*,mark size=1.4pt,samples at={-3,3}] {0.129};
\addplot+[ycomb,specred,very thick,mark=*,mark size=1.4pt,samples at={-4,4}] {0.034};
\node[sigblue] at (axis cs:0,0.30) {\footnotesize carrier $J_0$};
\end{axis}
\end{tikzpicture}
\caption{The line spectrum of a single-tone FM signal with modulation index $\beta=2$.
Sidebands sit at multiples of $\omega_m$ from the carrier, and their heights are the
Bessel values $J_n(\beta)$. The carrier line ($J_0$) is not even the tallest -- power has
spilled into the sidebands. Carson's rule \eqref{ch23:carson} captures where the comb has
effectively died out.}
\label{ch23:fig-fmspec}
\end{figure}

\begin{historybox}
Edwin Armstrong demonstrated wideband FM in 1933, against the prevailing belief (voiced
by no less than the mathematician John Carson) that frequency modulation could only
\emph{narrow} bandwidth, never improve fidelity. Armstrong's insight was the opposite of
intuition: by \emph{spending} bandwidth on large frequency deviation, FM trades it for a
dramatic gain in signal-to-noise ratio, because noise mostly perturbs amplitude, and an
FM receiver throws amplitude away. The legal and corporate battles that followed were as
bitter as the engineering was brilliant.
\end{historybox}

\begin{whybox}
\textbf{Why FM resists noise.} Most noise adds to the \emph{amplitude} of a received
wave; it nudges the arrow's length far more than its rotation rate. An FM receiver begins
by hard-limiting -- clipping the signal to constant amplitude -- which discards those
length perturbations entirely, keeping only the timing of the zero crossings, where the
message lives. AM has no such escape: its message \emph{is} the amplitude, so amplitude
noise lands directly on the recovered signal. This is the deep reason a wrist sensor that
encodes information in \emph{timing} (zero-crossing rate, frequency) can be more robust
than one that encodes it in \emph{level}.
\end{whybox}

\section{The analytic signal: a preview}
\label{ch23:analytic}

Both AM and FM forced us to talk about ``the envelope'' and ``the instantaneous
frequency'' of a real signal. But a real cosine is ambiguous: $A\cos(\omega t)$ is two
phasors spinning in opposite directions, and ``the'' rotation rate of a counter-rotating
pair is not well defined. To make envelope and instantaneous frequency precise, we keep
\emph{only one} of the two phasors. That one-sided object is the \emph{analytic signal}.

Given a real signal $s(t)$, its analytic signal is
\begin{equation}
\label{ch23:analytic-def}
s_a(t) = s(t) + \jj\,\hat{s}(t),
\end{equation}
where $\hat{s}(t)$ is the \emph{Hilbert transform} of $s$ -- the next chapter is devoted
to it. Geometrically, the Hilbert transform supplies the missing imaginary part so that
the negative-frequency phasor cancels and only the positive-frequency one remains. The
test case makes this vivid: for $s(t)=\cos\omega_c t$ we have $\hat s(t)=\sin\omega_c t$,
so
\begin{equation}
s_a(t) = \cos\omega_c t + \jj\sin\omega_c t = e^{\jj\omega_c t},
\end{equation}
a single clean phasor. In the frequency domain, forming $s_a$ deletes the negative-
frequency half of the spectrum and doubles the positive half.

Once we have the single arrow $s_a(t)=a(t)\,e^{\jj\phi(t)}$, the envelope and
instantaneous frequency we used informally above become exact:
\begin{equation}
\label{ch23:envphase}
a(t) = \abs{s_a(t)} \quad(\text{envelope}),
\qquad
\omega_i(t) = \frac{\dd}{\dd t}\arg s_a(t) \quad(\text{instantaneous frequency}).
\end{equation}

\begin{happeningbox}
The analytic signal is the bridge from real, two-sided, ambiguous waveforms to a single
unambiguous spinning arrow. Everything we are about to call ``I/Q'' is just the analytic
signal after we have slid it down to baseband: stop the fast spin, and what remains is a
slowly drifting complex number whose length is the envelope and whose angle is the phase.
The Hilbert chapter proves the cancellation; here we simply use the result.
\end{happeningbox}

\section{I/Q: the complex baseband}
\label{ch23:iq}

Here is the central construction of the chapter. Take a real bandpass signal -- a carrier
at $\omega_c$ carrying some message -- and multiply it by \emph{both} a cosine and a sine
at the carrier frequency. The two products are the \emph{in-phase} ($I$) and
\emph{quadrature} ($Q$) components:
\begin{equation}
\label{ch23:iq-mix}
I(t) = \text{LPF}\big\{\,s(t)\cdot 2\cos\omega_c t\,\big\},
\qquad
Q(t) = \text{LPF}\big\{\,s(t)\cdot (-2\sin\omega_c t)\,\big\},
\end{equation}
where LPF is a low-pass filter that throws away the copy that landed up at $2\omega_c$.
``Quadrature'' just means $90^\circ$ apart: cosine and sine are the same wave a quarter
cycle out of step. Bundle the two real numbers into one complex number, the
\emph{complex baseband} or \emph{I/Q signal}:
\begin{equation}
\label{ch23:iq-complex}
z(t) = I(t) + \jj\,Q(t).
\end{equation}

To see what $z$ is, write the bandpass signal in its general narrowband form,
$s(t) = a(t)\cos\!\big(\omega_c t + \phi(t)\big)$. Expanding and discarding the
$2\omega_c$ terms with the LPF gives the clean result
\begin{equation}
\label{ch23:iq-result}
I(t) = a(t)\cos\phi(t),\qquad Q(t)=a(t)\sin\phi(t),
\qquad\Longrightarrow\qquad
z(t) = a(t)\,e^{\jj\phi(t)}.
\end{equation}
So $z(t)$ is exactly the analytic signal with the carrier spin removed: its
\emph{magnitude} is the envelope and its \emph{angle} is the message phase. The single
complex number per sample carries both pieces of information that a single real number
cannot.

\begin{figure}[H]
\centering
\begin{tikzpicture}[scale=1.0]
% complex plane
\draw[->,gray] (-2.2,0)--(2.6,0) node[right]{$I$};
\draw[->,gray] (0,-2.2)--(0,2.6) node[above]{$Q$};
% the I/Q arrow
\def\ang{52}
\draw[sigblue,very thick,-{Stealth[length=3mm]}] (0,0)--(\ang:2.0);
\node[sigblue] at (\ang:2.3) {$z(t)$};
% projections
\draw[specred,dashed] (\ang:2.0)--(\ang:2.0 |- 0,0);
\draw[specred,dashed] (\ang:2.0)--(0,0 |- \ang:2.0);
\node[specred] at (0.62,-0.32) {$I$};
\node[specred] at (-0.34,1.3) {$Q$};
% angle and radius labels
\draw[targetgreen,thick] (0.7,0) arc (0:\ang:0.7);
\node[targetgreen] at (0.95,0.42) {$\phi$};
\node[sigblue] at (1.0,1.55) {$a=|z|$};
% spin hint
\draw[gray,->] (\ang+12:2.25) arc (\ang+12:\ang+42:2.25);
\node[gray] at (\ang+58:2.4) {\footnotesize spin $\dot\phi=\omega_i-\omega_c$};
\end{tikzpicture}
\caption{The I/Q sample as an arrow in the complex plane. Its horizontal projection is
$I$, its vertical projection is $Q$, its length is the envelope $a(t)$, and its angle is
the phase $\phi(t)$. After the carrier is removed, the arrow drifts slowly; its residual
spin rate is the message's instantaneous frequency offset from the carrier.}
\label{ch23:fig-iqplane}
\end{figure}

\begin{intuitionbox}
Imagine standing on a carousel turning at exactly the carrier rate $\omega_c$, watching
the received phasor. In that rotating frame the fast carrier spin is cancelled, and the
phasor appears nearly still -- drifting only as the \emph{message} changes it. The I/Q
mixer in \eqref{ch23:iq-mix} is precisely the act of stepping onto that carousel.
$\cos\omega_c t$ measures the part of the arrow pointing ``forward'' along your motion,
$\sin\omega_c t$ the part pointing ``sideways.'' Together they pin down the arrow's tip
in the rotating frame, which is exactly $z(t)$.
\end{intuitionbox}

\begin{formalbox}
The two carriers $\cos\omega_c t$ and $\sin\omega_c t$ are \emph{orthogonal} over any
interval that is an integer number of carrier periods:
$\int_0^{T}\cos(\omega_c t)\sin(\omega_c t)\,\dd t = 0$. That orthogonality is why $I$ and
$Q$ are independent channels -- projecting onto one tells you nothing about the other.
In the inner-product language of the linear-algebra companion, $\{\cos\omega_c t,\,
\sin\omega_c t\}$ is an orthogonal basis for the two-dimensional space of sinusoids at
$\omega_c$, and $(I,Q)$ are simply the coordinates of $s$ in that basis. This is the same
geometry as the real/imaginary split of a complex number.
\end{formalbox}

\subsection{Reconstructing the real signal from I/Q}
\label{ch23:iq-up}

The construction is reversible. Given the complex baseband $z=I+\jj Q$, you rebuild the
real passband signal by multiplying back up by the carrier and taking the real part:
\begin{equation}
\label{ch23:upconvert}
s(t) = \Re\!\big\{ z(t)\,e^{\jj\omega_c t}\big\}
     = I(t)\cos\omega_c t - Q(t)\sin\omega_c t.
\end{equation}
The two-multiplier structure on the right is the \emph{quadrature modulator} (or ``I/Q
upconverter''), and the pair (mixer down, mixer up) is the backbone of every software-
defined radio transmitter and receiver. Different choices of how you fill $I(t)$ and
$Q(t)$ give every modulation scheme at once: put the message in $\abs{z}$ and you have
AM; put it in $\arg z$ and you have FM/PM; put discrete points in the $(I,Q)$ plane and
you have the digital constellations (QPSK, 16-QAM) behind Wi-Fi and LTE.

\begin{figure}[H]
\centering
\begin{tikzpicture}[font=\small,>=Stealth]
% blocks
\node (in) at (0,0) {$s(t)$};
\node[draw,circle,inner sep=1.5pt] (m1) at (2.2,0.9) {$\times$};
\node[draw,circle,inner sep=1.5pt] (m2) at (2.2,-0.9) {$\times$};
\node[draw,rounded corners] (l1) at (4.0,0.9) {LPF};
\node[draw,rounded corners] (l2) at (4.0,-0.9) {LPF};
\node (Iout) at (5.7,0.9) {$I(t)$};
\node (Qout) at (5.7,-0.9) {$Q(t)$};
% wires
\draw[->] (in.east) -- (0.9,0) |- (m1.west);
\draw[->] (0.9,0) |- (m2.west);
\draw[->] (m1) -- (l1);
\draw[->] (m2) -- (l2);
\draw[->] (l1) -- (Iout);
\draw[->] (l2) -- (Qout);
% local oscillator labels
\node[sigblue] at (2.2,2.0) {$2\cos\omega_c t$};
\draw[->,sigblue] (2.2,1.75) -- (m1.north);
\node[sigblue] at (2.2,-2.0) {$-2\sin\omega_c t$};
\draw[->,sigblue] (2.2,-1.75) -- (m2.south);
\end{tikzpicture}
\caption{The I/Q downconverter (quadrature demodulator). The incoming real signal is
multiplied by a cosine and a sine from the local oscillator; each product is low-pass
filtered to drop the image at $2\omega_c$, yielding the baseband pair $I(t),Q(t)$. Run
the diagram backward, summing instead of splitting, to get the upconverter of
\eqref{ch23:upconvert}.}
\label{ch23:fig-block}
\end{figure}

\section{Why I/Q is everywhere}
\label{ch23:why-iq}

It is worth stating plainly why this complex-baseband representation took over radio and
sensing, because the reasons are practical, not aesthetic.

\textbf{1. Negative frequencies become real.} A real signal's spectrum is symmetric:
the negative-frequency half is a mirror of the positive half and carries no new
information. A \emph{complex} baseband signal has an \emph{asymmetric} spectrum -- it can
have energy at $+3\,$kHz and none at $-3\,$kHz. That asymmetry is exactly what lets I/Q
distinguish a tone \emph{above} the local oscillator from one \emph{below} it. A single
real mixer cannot tell those apart (both fold to the same baseband frequency); the I/Q
pair can, because the sign of $Q$ relative to $I$ encodes which side of the carrier the
energy came from.

\textbf{2. Half the sampling rate.} A real bandpass signal of bandwidth $B$ centered at a
high $\omega_c$ would need a fast ADC to sample directly. After I/Q downconversion the
information sits in two slow baseband channels, each centered on zero. Complex sampling at
rate $f_s$ covers a full bandwidth $f_s$, versus only $f_s/2$ for real sampling, so the
$(I,Q)$ pair captures the same band at half the per-channel rate of a direct real ADC.
The hardware shrinks dramatically.

\textbf{3. One representation for every modulation.} As noted, AM, FM, PM, and all
digital constellations are just trajectories of the same arrow $z(t)$ in the $(I,Q)$
plane. A software-defined radio implements them all by changing arithmetic on $I+\jj Q$,
not by changing circuits.

\textbf{4. The math is the analytic signal.} I/Q \emph{is} the analytic signal shifted to
baseband, so all the clean facts -- envelope is $\abs{z}$, instantaneous frequency is the
derivative of $\arg z$ -- apply directly, with no $2\omega_c$ junk to filter mentally.

\begin{figure}[H]
\centering
\begin{tikzpicture}[scale=1.05]
\draw[->,gray] (-2.1,0)--(2.3,0) node[right]{$I$};
\draw[->,gray] (0,-2.1)--(0,2.3) node[above]{$Q$};
% four QPSK points at (+-1,+-1)/sqrt2 ~ 0.707
\foreach \x/\y in {0.707/0.707, -0.707/0.707, -0.707/-0.707, 0.707/-0.707}{
  \filldraw[sigblue] (\x,\y) circle (2.2pt);
  \draw[sigblue,thick,-{Stealth[length=2mm]}] (0,0)--(\x,\y);
}
\node[sigblue] at (1.15,1.0) {\footnotesize $(+,+)$};
\node[sigblue] at (-1.15,1.0) {\footnotesize $(-,+)$};
\node[sigblue] at (-1.15,-1.0) {\footnotesize $(-,-)$};
\node[sigblue] at (1.15,-1.0) {\footnotesize $(+,-)$};
% a faint rotation showing frequency offset
\draw[specred,->,thick] (1.35:1.0) arc (1.35:35:1.0);
\node[specred] at (40:1.25) {\footnotesize $\Delta\omega\,t$};
\end{tikzpicture}
\caption{A QPSK constellation: each transmitted symbol is one of four fixed points in the
$(I,Q)$ plane, $(\pm1,\pm1)/\sqrt2$. A receiver local-oscillator frequency offset
$\Delta\omega$ makes the whole constellation rotate at rate $\Delta\omega$ (red arc) --
the same spinning that fades coherent DSB-SC in and out, seen as a picture.}
\label{ch23:fig-qpsk}
\end{figure}

\begin{applicationbox}
\textbf{I/Q in sensing, not just radio.} Doppler radar and ultrasound velocimetry read
target \emph{velocity} from the slow rotation rate of $z(t)$: a moving target shifts the
echo's carrier, so the baseband arrow spins at the Doppler frequency, and the
\emph{sign} of its spin (only visible with I/Q) tells you approaching versus receding.
The same trick powers $60\,$GHz gesture radars (e.g.\ Google's Soli) that watch a hand's
micro-motions as patterns in the I/Q time series. MRI scanners record the spin signal
directly as quadrature $I/Q$. Even a humble capacitive touch controller demodulates an
excitation tone into I and Q to separate ``how much capacitance'' (in phase) from
``how much loss'' (quadrature).
\end{applicationbox}

\begin{goalbox}
Our wrist controller does not transmit radio, but I/Q thinking pays off twice. First,
several MEMS sensor front ends are \emph{themselves} I/Q systems under the hood: a MEMS
gyroscope drives a proving mass at a fixed resonance and reads Coriolis-induced motion by
\emph{synchronous (I/Q) demodulation} against that drive tone -- the very lock-in scheme
of \eqref{ch23:iq-mix}. Understanding I/Q demystifies why a gyro has a ``drive
frequency'' in its datasheet and why temperature shifts of that frequency cause bias
drift (Chapter on calibration).

Second, we can borrow the technique deliberately. Suppose we attach a small piezo or
electrode that injects a steady tone into the forearm and listen for how a clench changes
it (a bioimpedance or force-myography idea). Demodulating that tone into $I$ and $Q$
gives us a clean, slowly varying two-number feature -- magnitude and phase of the
tissue's response -- far easier to threshold into a gesture than the raw high-frequency
waveform. \emph{The limit to keep in mind:} this yields a coarse ``how hard is the wrist flexing''
scalar, not individual finger joint angles. Continuous per-finger pose still requires
sensors on the fingers (instrumented rings) or rich multi-channel EMG/FMG arrays; a
single wrist tone, however cleverly demodulated, cannot recover what it never measured.
\end{goalbox}

\begin{buildbox}
\textbf{Software I/Q demodulation in a few lines.} If you sample a real signal at $f_s$
and know the carrier $f_c$, you can do the I/Q split entirely in code -- no analog mixer
required. Multiply by the complex carrier and low-pass filter:
\begin{lstlisting}[language=Python,caption={Numerical I/Q downconversion of a real bandpass signal.},label={ch23:lst-iq}]
import numpy as np
from scipy.signal import butter, filtfilt

def iq_demod(s, fs, fc, bw):
    n  = np.arange(len(s))
    lo = np.exp(-1j * 2*np.pi * fc * n / fs)   # carousel: spin backward at fc
    z  = s * lo                                # shift carrier copy down to 0 Hz
    b, a = butter(4, bw / (fs/2))              # low-pass to drop the 2*fc image
    I = filtfilt(b, a, z.real)
    Q = filtfilt(b, a, z.imag)
    return I + 1j*Q                            # complex baseband

# envelope and instantaneous frequency from the I/Q arrow:
# env   = np.abs(z)
# instf = np.diff(np.unwrap(np.angle(z))) * fs / (2*np.pi) + fc
\end{lstlisting}
On a microcontroller you would replace \texttt{filtfilt} with a small causal IIR or a
moving average, and precompute the sine/cosine tables. The structure is identical: spin
backward by the carrier, then smooth.
\end{buildbox}

\section{Demodulation: getting the message back}
\label{ch23:demod}

Demodulation is just \eqref{ch23:cosmod} run in reverse: multiply the received signal by a
carrier again to slide one of the sidebands back to the origin, then low-pass filter to
keep it and discard the copy that flew up to $2\omega_c$. There are two flavors,
distinguished by whether the receiver knows the carrier's phase.

\subsection{Coherent (synchronous) demodulation}
\label{ch23:coherent}

If the receiver has a local oscillator \emph{locked} to the carrier's frequency and
phase, recovery is exact. For DSB-SC, $s(t)=m(t)\cos\omega_c t$, multiply again:
\begin{equation}
\label{ch23:coh}
s(t)\cdot 2\cos\omega_c t = 2 m(t)\cos^2\omega_c t
 = m(t) + m(t)\cos 2\omega_c t .
\end{equation}
The first term is the message back at baseband; the second is a copy at $2\omega_c$ that
the low-pass filter removes. This is the I/Q downconverter of \eqref{ch23:iq-mix} with
$\phi=0$, and it is the gold standard: coherent detection is what gives digital radios
their noise performance, and it is mandatory for any phase-bearing scheme (FM, PM, QAM).

\begin{trapbox}
\textbf{Phase error costs you.} If the local oscillator is off the carrier phase by a
fixed $\Delta\phi$, coherent DSB-SC recovers $m(t)\cos\Delta\phi$ -- attenuated, and dead
($\cos 90^\circ=0$) if you are a quarter cycle off. A \emph{frequency} error
$\Delta\omega$ is worse: the recovered message gets multiplied by the slowly rotating
$\cos(\Delta\omega\, t)$, fading in and out. In I/Q terms, a frequency offset makes the
whole constellation \emph{spin}. Real receivers spend real effort on a carrier-recovery
loop (a phase-locked loop, or a digital equivalent) to drive $\Delta\phi,\Delta\omega\to
0$.
\end{trapbox}

\subsection{Envelope detection (noncoherent)}
\label{ch23:envdet}

For full AM \eqref{ch23:am-eq}, where the message rides in a strictly positive envelope,
you do not need the carrier phase at all. Just follow the peaks: rectify (take
$\abs{\cdot}$), then low-pass to smooth out the carrier ripple. The output tracks
$A[1+\mu m(t)]$, and a DC-blocking capacitor strips the constant to leave $m(t)$. This is
the diode-and-capacitor detector of the crystal radio, and the geometric statement of
``follow the peaks'' is simply ``read off $\abs{z(t)}$,'' the length of the I/Q arrow.

\begin{figure}[H]
\centering
\begin{tikzpicture}
\begin{axis}[width=13cm,height=4.3cm,
  xlabel={time $t$}, ylabel={amplitude},
  xmin=0,xmax=6.3,ymin=-0.1,ymax=2.1,
  ytick=\empty, xtick=\empty, axis lines=left,
  title={Envelope detection: rectify, then smooth, to read off $|z(t)|$}]
% rectified AM (positive half)
\addplot[gray,thin,domain=0:6.3,samples=200]
   {abs((1+0.7*cos(deg(x)))*cos(deg(13*x)))};
% recovered envelope
\addplot[specred,very thick,domain=0:6.3,samples=120] {1+0.7*cos(deg(x))};
\node[specred] at (axis cs:3.1,0.4) {\footnotesize recovered $A[1+\mu m(t)]$};
\end{axis}
\end{tikzpicture}
\caption{Noncoherent envelope detection. The gray trace is the rectified AM waveform; the
smoothing filter rides its peaks to recover the envelope (red), which is the message plus
a constant. No carrier phase knowledge is needed -- only that the envelope never crossed
zero.}
\label{ch23:fig-envdet}
\end{figure}

\subsection{FM demodulation}
\label{ch23:fmdemod}

To demodulate FM, recover the instantaneous frequency, which by \eqref{ch23:envphase} is
the derivative of the phase of the analytic (I/Q) signal. The cleanest digital recipe is:
downconvert to I/Q, then differentiate the unwrapped angle,
\begin{equation}
\label{ch23:fmdemod-eq}
\hat m(t) \;\propto\; \omega_i(t)-\omega_c \;=\; \frac{\dd}{\dd t}\arg z(t).
\end{equation}
A neat trick avoids the explicit arctangent: since
$\frac{\dd}{\dd t}\arg z = \Im\{\dot z\,\bar z\}/\abs{z}^2$, one differentiates $z$,
multiplies by the conjugate, and takes the imaginary part. Hardware FM detectors
approximate the same derivative with a slope filter or a phase-locked loop. Either way,
the message is read from \emph{how fast the I/Q arrow is spinning}, exactly the picture in
Figure~\ref{ch23:fig-iqplane}.

\begin{connectionbox}
The ``differentiate the angle'' recipe leans on the chain rule and on the fact (calculus
companion) that $\frac{\dd}{\dd t}\arg z = \Im(\dot z/z)$. The ``multiply by the
conjugate'' form is the same expression cleared of its denominator's phase -- a small
piece of complex algebra that saves an expensive $\arctan$ and an unwrap on an embedded
chip. It reappears in pitch trackers and in phase-vocoder time stretching in the audio
chapters.
\end{connectionbox}

\section{A worked numeric example}
\label{ch23:worked}

Let us push numbers through to make the spectra concrete. Take a message tone at
$f_m = 1\,$kHz, $m(t)=\cos(2\pi\cdot 1000\,t)$, and a carrier at $f_c = 20\,$kHz.

\textbf{DSB-SC.} The product $s(t)=m(t)\cos(2\pi f_c t)$ has lines at
$f_c\pm f_m = 19\,$kHz and $21\,$kHz (and their negatives). No energy at $20\,$kHz: the
carrier is suppressed. Bandwidth occupied: $2f_m = 2\,$kHz.

\textbf{Full AM, $\mu=0.5$.} $s(t)=A[1+0.5\cos(2\pi f_m t)]\cos(2\pi f_c t)$ adds a
carrier line at $20\,$kHz of relative height $1$ plus sidebands at $19$ and $21\,$kHz of
relative height $0.25$ each. The fraction of power in the (informative) sidebands is
$\frac{\mu^2/2}{1+\mu^2/2} = \frac{0.125}{1.125}\approx 11\%$ -- a vivid reminder of why
suppressing the carrier is attractive.

\textbf{I/Q sanity check.} Downconvert the DSB-SC signal by multiplying by
$e^{-\jj 2\pi f_c t}$ and low-pass filtering. The $21\,$kHz line moves to $+1\,$kHz, the
$19\,$kHz line to $-1\,$kHz; the negative-frequency images move to $-41$ and $-39\,$kHz
and are filtered away. The surviving baseband $z(t)=\tfrac12 e^{\jj 2\pi f_m t}+\tfrac12
e^{-\jj 2\pi f_m t}=\cos(2\pi f_m t)=m(t)$ -- real, because DSB-SC had a symmetric
spectrum about the carrier. Had we sent only the upper sideband (SSB), $z$ would be the
single complex tone $\tfrac12 e^{\jj 2\pi f_m t}$, asymmetric, and the imaginary part of
$z$ would be nonzero -- the I/Q pair \emph{seeing} the one-sidedness that a real mixer
would miss.

\begin{figure}[H]
\centering
\begin{tikzpicture}
\begin{axis}[width=13cm,height=4.2cm,
  xlabel={frequency (kHz)}, ylabel={relative height},
  xmin=17.5,xmax=22.5,ymin=0,ymax=1.15,
  xtick={19,20,21}, ytick=\empty,
  axis lines=left, tick align=outside,
  title={Spectrum of full AM with $\mu=0.5$: tall carrier, short sidebands}]
\addplot+[ycomb,sigblue,very thick,mark=*,mark size=1.6pt,samples at={20}] {1.0};
\addplot+[ycomb,specred,very thick,mark=*,mark size=1.6pt,samples at={19,21}] {0.25};
\node[sigblue] at (axis cs:20,1.07) {\footnotesize carrier};
\node[specred] at (axis cs:18.55,0.36) {\footnotesize LSB};
\node[specred] at (axis cs:21.5,0.36) {\footnotesize USB};
\end{axis}
\end{tikzpicture}
\caption{Positive-frequency line spectrum of the worked full-AM example. The carrier at
$20\,$kHz dwarfs the two sidebands at $19$ and $21\,$kHz, where all the information lives
-- the power inefficiency that motivates suppressed-carrier and single-sideband schemes.}
\label{ch23:fig-amspec}
\end{figure}

\begin{labbox}
Generate the full-AM signal of the worked example at a sample rate of $200\,$kHz, take its
FFT, and confirm the three lines at $19,20,21\,$kHz with the $4{:}1$ height ratio between
carrier and sidebands. Then I/Q-downconvert with \texttt{iq\_demod} from
Listing~\ref{ch23:lst-iq} and verify that $\abs{z(t)}$ reproduces $1+0.5\cos(2\pi\cdot
1000\,t)$. Finally, change $\mu$ to $1.5$ and watch envelope detection fail where the
envelope folds -- the overmodulation trap, seen rather than merely warned about.
\end{labbox}

\section{Summary}
\label{ch23:summary}

Modulation is the single idea that multiplying by a sinusoid slides a spectrum along the
frequency axis -- a corollary of the convolution theorem. AM writes the message onto the
carrier's height (easy to detect, power-inefficient, vulnerable to amplitude noise); FM
writes it onto the carrier's rotation rate (constant amplitude, wider bandwidth by
Carson's rule, robust to amplitude noise). The analytic signal removes the ambiguity of
real cosines by keeping one phasor, making ``envelope'' and ``instantaneous frequency''
precise. Sliding the analytic signal down to baseband gives the I/Q pair $z=I+\jj Q$, a
single complex number per sample whose length is the envelope and whose angle is the
phase. I/Q is everywhere -- it distinguishes positive from negative frequencies, halves
sampling rates, unifies every modulation scheme, and underlies sensors from Doppler radar
to the MEMS gyro on our own wrist. Demodulation is just modulation in reverse: coherent
detection when phase is known, envelope detection when it is not, and angle-
differentiation for FM. The next chapter makes the analytic signal rigorous by building
the Hilbert transform that supplies its imaginary half.

\section{Exercises}
\label{ch23:exercises}

\begin{exercise}[The shift, by hand]
\label{ch23:ex-shift}
Starting from the analysis integral, prove the modulation property
\eqref{ch23:cosmod} for $\cos\omega_c t$ by writing the cosine as
$\tfrac12(e^{\jj\omega_c t}+e^{-\jj\omega_c t})$ and applying \eqref{ch23:expshift} to
each term. Identify which copy is the upper sideband and which is the lower.
\end{exercise}

\begin{exercise}[Sidebands of two tones]
\label{ch23:ex-twotone}
A message $m(t)=\cos(2\pi\cdot 800\,t)+\tfrac12\cos(2\pi\cdot 2000\,t)$ amplitude-
modulates a $50\,$kHz carrier in full-AM form with $\mu=0.6$. List every spectral line
(in kHz) on the positive-frequency axis and its relative height. What total bandwidth does
the signal occupy?
\end{exercise}

\begin{exercise}[Power in the sidebands]
\label{ch23:ex-power}
For full AM with a single sinusoidal message and modulation index $\mu$, show that the
fraction of total transmitted power carried by the two sidebands is
$\dfrac{\mu^2/2}{1+\mu^2/2}$. Evaluate it at $\mu=0.5$ and $\mu=1$, and comment on why
broadcasters nonetheless keep $\mu$ below $1$.
\end{exercise}

\begin{exercise}[Instantaneous frequency]
\label{ch23:ex-instf}
A signal is $s(t)=\cos(\omega_c t + \beta\sin\omega_m t)$. Compute its instantaneous
frequency $\omega_i(t)$ using \eqref{ch23:instfreq}. What are the maximum and minimum
instantaneous frequencies, and how do they relate to the peak deviation $\Delta f$ in
Carson's rule \eqref{ch23:carson}?
\end{exercise}

\begin{exercise}[Carson's rule]
\label{ch23:ex-carson}
An FM broadcast station has peak frequency deviation $\Delta f = 75\,$kHz and a highest
audio frequency $f_m = 15\,$kHz. Estimate the occupied bandwidth from Carson's rule
\eqref{ch23:carson}. How does this compare to the $200\,$kHz channel spacing used in FM
broadcasting, and why is the guard band reasonable?
\end{exercise}

\begin{exercise}[I/Q of a known signal]
\label{ch23:ex-iq}
Take $s(t)=\cos\!\big((\omega_c+\omega_m)t\big)$ -- a single upper-sideband tone. Apply
the I/Q downconversion \eqref{ch23:iq-mix} (multiply by $2\cos\omega_c t$ and
$-2\sin\omega_c t$, low-pass) and show that $z(t)=e^{\jj\omega_m t}$. Now repeat for a
\emph{lower}-sideband tone $\cos((\omega_c-\omega_m)t)$ and show $z(t)=e^{-\jj\omega_m
t}$. Explain how the \emph{sign} of $z$'s rotation distinguishes the two, and why a single
real mixer cannot.
\end{exercise}

\begin{exercise}[Phase-error penalty]
\label{ch23:ex-phaseerr}
A coherent DSB-SC receiver multiplies $s(t)=m(t)\cos\omega_c t$ by $2\cos(\omega_c
t+\Delta\phi)$ and low-pass filters. Show that the recovered baseband is
$m(t)\cos\Delta\phi$. At what $\Delta\phi$ does the output vanish, and what does this
imply about the need for carrier recovery? Then redo the calculation with a frequency
error $\Delta\omega$ and describe the result.
\end{exercise}

\begin{exercise}[Envelope versus coherent]
\label{ch23:ex-envcoh}
Explain in two or three sentences why an envelope detector works for full AM ($\mu\le 1$)
but fails for DSB-SC. Then explain why DSB-SC \emph{can} be recovered coherently. Tie both
answers to the geometry of the I/Q arrow $z(t)$ and whether its length alone, or its full
position, carries the message.
\end{exercise}

\begin{exercise}[FM discriminator without arctan]
\label{ch23:ex-fmdisc}
Verify the identity $\dfrac{\dd}{\dd t}\arg z(t) = \dfrac{\Im\{\dot z(t)\,\overline{z(t)}\}}
{\abs{z(t)}^2}$ for a complex signal $z(t)=a(t)e^{\jj\phi(t)}$ with $a(t)>0$. Explain why
the numerator alone (without dividing by $\abs{z}^2$) is an adequate FM discriminator when
the envelope is constant, and why this is attractive on a microcontroller.
\end{exercise}

\begin{exercise}[Numerical experiment]
\label{ch23:ex-numeric}
Using Listing~\ref{ch23:lst-iq}, synthesize an FM signal whose message is a $200\,$Hz
sine, carrier $10\,$kHz, deviation $1\,$kHz, sampled at $100\,$kHz. Downconvert to I/Q,
recover the instantaneous frequency by differentiating the unwrapped angle, and plot it
against the true message. Where does the recovered message degrade, and what filter
bandwidth in \texttt{iq\_demod} gives the cleanest result?
\end{exercise}

\begin{exercise}[Lock-in / synchronous detection]
\label{ch23:ex-lockin}
A tiny sensor signal $\eps\cos(\omega_c t+\phi)$ is buried in broadband noise. You
multiply the noisy measurement by $\cos\omega_c t$ and by $\sin\omega_c t$ and low-pass
filter each. Show that the two outputs recover $\tfrac{\eps}{2}\cos\phi$ and
$\tfrac{\eps}{2}\sin\phi$, so that magnitude and phase of the buried tone are obtained
while out-of-band noise is rejected. Relate this to the MEMS-gyro readout described in the
``Toward the wrist tracker'' box.
\end{exercise}

\begin{exercise}[Constellation thinking]
\label{ch23:ex-constellation}
In QPSK, the transmitter sets $(I,Q)$ to one of the four points $(\pm 1,\pm 1)/\sqrt2$
once per symbol (Figure~\ref{ch23:fig-qpsk}). Sketch these on the $(I,Q)$ plane. If the receiver's local oscillator has
a slow frequency offset $\Delta\omega$, describe what happens to the received constellation
over time, and explain why this is the same phenomenon as the fading $\cos(\Delta\omega\,
t)$ term in coherent DSB-SC demodulation.
\end{exercise}
